\scanpage{54}

% manuscript line(s) 1–2
Si $y=(\lambda,\unsure{\zeta|\xi})\in\Gamma P_1$ \deleted{\unsure{alors}} \inserted{i.e. $\lambda=1$,} l'équation $Q(y)=0$ équivaut bien à $(\zeta^2\deleted{\unsure{+}}\unsure{-|+}\delta)/4=0$ i.e.
% manuscript line(s) 3
\[
\zeta^2=\delta
\]
% manuscript line(s) 4–7
(en plus on a $\zeta_0=T_0$, sous-entendu), cela prouve que $Q'(\lambda,\zeta)=(\zeta^2\deleted{\unsure{+}}-\lambda^2\delta)/4$ est prop. à $Q$, et pour voir que le facteur de prop. est $1$, il suffit de faire \deleted{$y=(\lambda,\zeta)=T$}
% manuscript line(s) 8
\[
y=T,\qquad\text{i.e.}\quad y=(2,0)
\]
% manuscript line(s) 9
on a alors
% manuscript line(s) 9
\[
\deleted{Q(T)=T^2-2}
\]
% manuscript line(s) 10–11
\[
\begin{aligned}Q(y)&=y^2-Ty\psi(y)+N\psi(y)^2=T^2-2T^2+4N={}\\ &4N-T^2=\deleted{\unsure{+}}-\delta\end{aligned}
\]
% manuscript line(s) 12
et $Q'(y)=(0+4\delta)/\deleted{\unsure{8}}4=\unsure{-|+}\delta$ \qquad OK.
% manuscript line(s) 13
\par\centerline{\rule{1.1cm}{.3pt}}

% manuscript line(s) 14
Comparaison des torseurs sous $\mu_{2}$ et sous $(\mathbb Z/2\mathbb Z)_X$.
% manuscript line(s) 15
L'hom can.
% manuscript line(s) 16
\[
(\mathbb Z/2\mathbb Z)_X\longrightarrow\mu_{2\,X}
\]
% manuscript line(s) 17
définit, par torseur associé, un $\otimes$-foncteur
% manuscript line(s) 18–21
\diagram{\begin{tikzcd}[ampersand replacement=\&,column sep=small,row sep=large] \underline{\mathrm{Tors}}\ (\mathbb Z/2\mathbb Z)\arrow[r,"\mathrm{Ass}","\deleted{\unsure{\varphi}}"']\arrow[d,"\unsure{\wr}","\deleted{\unsure{\mathcal A}}"']\&\mathrm{Tors}(\mu_2\deleted{\unsure{,X}})\arrow[r,"\mathrm{Oub}"]\arrow[d,"\unsure{\alpha|\sim|\wr}"]\&{\big[\mathrm{Rev\,quad}(X)\big]}\arrow[d]\\ {\substack{\underline{\text{triples}}\ (L,\underset{\substack{\in\\\Gamma L^{\otimes2}}}{\delta},\underset{\substack{\in\\\Gamma L_0}}{T_0})\\(\underline{\delta\ \mathrm{inv.}})}}\arrow[r,"{\substack{\mathrm{Ass}'\\?}}"]\&{\substack{\underline{\text{couples}}\ (L,\underset{\substack{\in\\\Gamma L^{\otimes2}}}{\Theta})\\(\underline{\Theta\ \mathrm{inv}})}}\arrow[r,"{\substack{\mathrm{Oub}'\\?}}"]\&{\Big[\,\substack{\underline{\mathrm{Triples}}\ (L,\underset{\substack{\in\\\Gamma L^{\otimes2}}}{\delta},\underset{\substack{\in\\\Gamma L_0}}{T_0})\\(\delta\text{ pas néc\textsuperscript{t} inv.})}}\end{tikzcd}}
% manuscript line(s) 22–24
Les flèches verticales sont des équiv. si $2$ rég. dans $X$ (et celles \inserted{du milieu sans condition sur $X$}) Il faut déterminer les deux foncteurs horizontaux, notés ``?''. Le premier foncteur est sûrement
% manuscript line(s) 25
\[
\mathrm{Ass}'(L,\delta,T_0)\unsure{\longrightarrow|=}(L,\delta)
\]
% archive sheet number: 53
